\chapter*{Conventions and notation}
\addcontentsline{toc}{chapter}{Conventions and notation}
\markboth{Conventions and notation}{Conventions and notation}

Unless a chapter says otherwise, \(\Bbbk\) is an algebraically closed
field. Chapters~8--10 work over \(\CC\) whenever classical topology,
equivariant cohomology, or symplectic language is used.

\begin{center}
\renewcommand{\arraystretch}{1.18}
\begin{tabularx}{\textwidth}{@{}>{\raggedright\arraybackslash}p{0.26\textwidth}X@{}}
\toprule
\textbf{Notation} & \textbf{Meaning} \\
\midrule
\(T=T_N\) & an algebraic torus \\
\(M=X^*(T)\) & the character lattice \\
\(N=X_*(T)=\Hom_{\ZZ}(M,\ZZ)\) & the lattice of one-parameter subgroups \\
\(\character{m}\), \(\oneps{u}\) & the character associated to \(m\in M\), and the one-parameter subgroup associated to \(u\in N\) \\
\(\pair{m}{u}\) & the perfect pairing \(M\times N\to\ZZ\), characterized by \(\character{m}(\oneps{u}(t))=t^{\pair{m}{u}}\) \\
\(M_{\RR},N_{\RR}\) & the real scalar extensions of the two lattices \\
\(\sigma^\vee\) & the dual cone of \(\sigma\subseteq N_{\RR}\) \\
\(U_\sigma\) & the affine toric variety \(\Spec\Bbbk[\sigma^\vee\cap M]\) \\
\(\Sigma\), \(X_\Sigma\) & a fan in \(N_{\RR}\), and its toric variety \\
\(\Sigma(1)\) & the set of rays of \(\Sigma\) \\
\(u_\rho\) & the primitive lattice generator of a ray \(\rho\in\Sigma(1)\) \\
\(O(\sigma)\), \(V(\sigma)\) & the orbit associated to \(\sigma\), and its closure \\
\(D_\rho\) & the invariant prime divisor corresponding to \(\rho\) \\
\(\Cl(X)\), \(\Pic(X)\) & the divisor class group and Picard group \\
\(P_D\subseteq M_{\RR}\) & the divisor polyhedron \(\{m:\pair{m}{u_\rho}\geq-a_\rho\}\) for \(D=\sum_\rho a_\rho D_\rho\) \\
\(\psi_D\) & the support function normalized by \(\psi_D(u_\rho)=-a_\rho\) \\
\(\overline{\NE}(X)\), \(\Nef(X)\) & the Mori cone of curves and the nef cone of divisors \\
\bottomrule
\end{tabularx}
\end{center}

We write \(\Cone{u_1,\ldots,u_r}\) for the convex polyhedral cone
generated by the displayed vectors and \(\Conv\{m_1,\ldots,m_s\}\)
for their convex hull. The symbol \(u_\rho\) is reserved for a
primitive ray generator; letters such as \(u_i,v_i,w_i\) may be used
locally when an example comes with an ordered list of rays or weights.
